This project analyzed historical crude oil price data using a 12 month moving average to highlight trends within the data. By comparing these trends with major global economic events, we are able to observe how external factors such as oil supply shocks, financial crisis, and global pandemics influence oil prices.

The results show that oil prices are highly sensitive to disruptions in global supply and demand. The data used shows how government policies can cause oil production to increase in order to counteract the high demand. While also seeing how companies are willing to cut production in order to increase prices and cause the population to pay more at the pump. We can canlayzze human factors as well. Whenever crisis are observed we can see the connection of crisis to buying in excess of products. The use of this data is significant in conveying how events can influence price behavior and trend direction. 

Overall, this analysis demonstrates how applying a 12 month moving average to a dataset, and applying historical context. We can understand the economic system and the factors that play a role in it.
